\section{Computational evidence and independent solution audits}
\label{sec:experiments}

The results establish three practical properties of the interface. First,
a frozen uniform experiment yields accepted bound artifacts for 203 of 289
attempted models in separate replay, with 80 accepted bounds satisfying $0\le d\le10^{-2}$ for the signed
normalized reference difference defined below.
Those references are unverified, so proximity does not establish optimality.
Second, exact audits distinguish invalid supplied proof steps and infeasible
returned points from failures of sufficient nonlinear checks. Third, saved
proofs make checking independent of repeating numerical search, but their
storage and replay costs remain substantial. We examine capability, failure
mechanisms, and costs in that order. Appendix~\ref{app:experimental-accounting}
retains the complete protocol, production accounting, and separately versioned
repair results; none of those later checks replaces the frozen uniform outcome.

\subsection{Capability: checkable bounds and their strength}
\label{sec:replay-results}

The frozen repository catalogue contains 299 models selected for convexity and
discrete variables from MINLPLib~\citep{minlplib}. Seven selected files did not
load and three did not pass the then-used curvature screen. The remaining 289
names form the complete denominator of the historical and uniform campaigns. The supplement
lists all 299 names and the ten exclusions, including their recorded reasons.
The selection is a repository snapshot, not a claim about the current size or
classification of MINLPLib. We distribute the loaded Python model files with
MINLPLib's CC-BY~4.0 attribution. GAMS is the library's primary representation;
the certificates concern the exact loaded expression semantics of
Section~\ref{sec:model}, not an unverified equivalence of every converted model.

The uniform protocol gives numerical OA and each exact-SCIP attempt a
requested 30-second search budget, allows at most a default and a conservative
fallback attempt, and caps each complete instance process group at 360 seconds.
It uses six workers, followed by a separate execution of the same checker with
six workers and a 1,200-second per-record limit. The search budgets exclude
additional preprocessing, proof completion, and checking work covered by the
outer cap. All 289 names remain in the denominator, including failed generation.

Separate replay accepts 203 artifacts, rejects 19, and finds 67 without a
complete bundle (Table~\ref{tab:campaign-results}). The producer itself returns
198 accepted results. Four additional accepted artifacts survive outer
timeouts and one survives a worker error. Thus an unsuccessful process return
need not destroy a completed checkable proof. This distinction measures the
availability of justified bounds under the stated protocol, rather than only
the producer's return status. It does not establish coverage of all convex
MINLPs or superiority over another solver.

\begin{table}[tb]
\centering
\caption{Complete replay outcomes and exact reference comparisons. Both denominators are 289; generation outcomes are reported separately in the text. Historical replay also used external corroboration. Reference rows count only verified bounds; $d$ is the signed normalized difference from an unverified reference (Section~\ref{sec:reference-comparison}).}
\label{tab:campaign-results}
\begin{tabular}{lrr}
\toprule
Quantity & Historical archive & Uniform-run artifacts \\
\midrule
Verified & 188 & 203 \\
Near reference: $0\le d\le10^{-4}$ & 52 & 46 \\
Near reference: $0\le d\le10^{-2}$ & 81 & 80 \\
Negative reference difference & 23 & 25 \\
Rejected & 92 & 19 \\
Missing complete artifacts & 9 & 67 \\
Timed out during replay & 0 & 0 \\
Other replay errors & 0 & 0 \\
\bottomrule
\end{tabular}
\end{table}


\paragraph{Strength relative to recorded objectives.}
\label{sec:reference-comparison}
For a replayed bound $B$ in the original objective sense, a recorded library
objective $r$, and $s=1$ for minimization or $s=-1$ for maximization, we compute
\[
        d(B,r)=\frac{s(r-B)}{\max\{1,|r|\}}.
\]
The bound and the decimal reference string are converted to exact rationals
before subtraction. We report the number satisfying $0\le d\le 10^{-4}$ and
$0\le d\le 10^{-2}$, and report $d<0$ separately. This is a signed normalized
difference from an unverified reference, not an optimality gap. Even $d=0$
does not prove that a feasible nonlinear point exists at that value.

Among the 203 accepted uniform-run bounds, 46 satisfy the first threshold
and 80 the second; 25 differences are negative. The interface therefore
produces some bounds close to recorded objectives, while the uniform success
count alone does not imply tight bounds. Among the 188 accepted historical
bounds, 52 satisfy the first threshold and 81 satisfy the second. The
historical comparison has 23 negative differences; every magnitude is
below $6\cdot10^{-10}$. Such small differences can arise from rounded reference
strings. In particular, we do not convert their sign into a claim of an
incorrect solver result. The supplement records every exact comparison,
including the uniform campaign's results in Table~\ref{tab:campaign-results}.
\paragraph{Bounds versus completed optimality proofs.}
A close reference value supplies no primal witness. Exact original-model
feasibility checks do complete the bound for \texttt{clay0204m}, giving optimum
6545 (Section~\ref{sec:solver-audits}), and for the small quadratic example
below. These are positive examples of the stronger conclusion available when
the primal and bound evidence match; we do not infer a benchmark optimality
count from objective comparisons.

The core's solver-free replay includes a mixed-integer quadratic example,
$\min y$ subject to $(x-\tfrac12)^2\le y$, $x\in\{0,1\}$ and $0\le y\le2$.
Eleven checked cuts and 39 proof derivations establish the lower bound
$1/4$, attained by $(x,y)=(0,1/4)$. The \texttt{clay0204m} matching-point audit
and the targeted repaired examples exercise larger complete bundles. These
checks supplement the 161 targeted software tests; test counts are not a
formal proof of the implementation. Section~\ref{sec:formalization} specifies
the separate proof-assistant result and its boundary.

Beyond the uniform experiment, the retained historical and targeted repair
artifacts form a useful descriptive collection. The repair versions and their
separate denominators are specified in Appendix~\ref{app:repair-accounting}.
The descriptive merged catalogue contains verified bounds for 222 distinct loaded models. It chooses the strongest exact normalized lower bound among 405 accepted records from the historical, primary, producer-repair, and reporting-repair evidence, requiring identical model SHA-256 and objective sense for every comparison. Exact ties use a fixed cohort order. Each entry identifies its original-sense bound, complete proof, source cohort, and hashes. This union is a reusable bound collection; its coverage is not attributed to the uniform search budget.


\subsection{Failure mechanisms: what rejection establishes}
\label{sec:failure-mechanisms}

A rejection is informative only when tied to its checking obligation. The
uniform run's 67 admission errors occur at the curvature/domain check, before
master generation; they account for all 67 missing replay bundles. Its nineteen
rejections comprise eighteen proof-rule or proof-grammar failures and one
post-proof reporting failure. The latter concerns representation of an already
checked exact bound, not failure of its mathematical inferences. Separately
recorded repairs address that boundary and twelve producer normalization
failures without changing the proof, domain, curvature, propagation, or
master-matching rules (Appendix~\ref{app:repair-accounting}).

\paragraph{Invalid supplied discrete inferences.}
Three new completed proofs receive an external acceptance message but fail
local conditions of the internal rules. Two linear combinations mix opposite
inequality directions: in each, a strictly positive multiplier of an upper
branch row appears among lower-row terms. The offending rows have nonzero
coefficients and are not constant tautologies. In the third proof, an
\texttt{uns} step cites the same integer coordinate $u$ with branches
$u\le64$ and $u\ge66$. Integer activity 65 is uncovered. Other master rows might
in principle exclude that activity, but the submitted standalone disjunction
step does not establish such an exclusion. The supplement gives the exact
antecedents and multipliers. These observations concern the individual supplied
inferences and do not imply false final bounds.

These cases show why external acceptance does not discharge the internal
inference obligations. The exact row directions and uncovered integer activity
supply local witnesses that are stronger evidence than process status or a
small objective discrepancy.

The historical audit provides a broader check on the effect of enforcing all
layers of the contract. Of 269 historical records labeled accepted, 188
retain that verdict. The 81
revocations comprise 67 domain or curvature failures, 13 cut-verification
failures, and one discrete-proof failure. Eleven further discrete-proof
failures came from historical records already labeled as crashed, giving 92
current rejections in total. Nine records lack a completed proof. No record in
the complete historical replay timed out or mutated during checking.
These counts replace the historical acceptance claim; the archived old flags
are retained solely to make the change auditable.

The twelve historical failures that reach the discrete proof checker admit
small exact arithmetic audits. In eleven, the claimed upper row has exactly
the coefficients obtained by the supplied rational linear combination, but
its right-hand side is smaller than that combination permits. The exact
positive excess ranges from approximately $1.55\cdot10^{-12}$ to
$1.93\cdot10^{-10}$. An arbitrarily small excess still invalidates that
inference over the reals. In the remaining instance,
\texttt{smallinvDAXr5b200-220}, a positive multiplier uses an upper inequality
in a combination requiring lower inequalities; the multiplied directions are
incompatible. The portable supplement includes the twelve first-failure
steps, their referenced rows, and an independent \texttt{Fraction}-arithmetic
recombination script. Full completed proofs are in the bulk archive.
These are invalid justifications, not proofs that their final proposed bounds
are unattainable or false.

\paragraph{Conservative nonlinear failures.}
Nonlinear rejection has a different interpretation. For example, the first
historical \texttt{risk2bpb} cut exceeds the newly recomputed sufficient safe
intercept by approximately $1.17\cdot10^{-12}$. This establishes failure of
this certificate test, not global invalidity of the affine inequality.
A separately recorded repaired generation for \texttt{risk2bpb} passes all
checks and gives an original-sense bound of approximately $-55.876139395$.
The library string $-55.8761394$ lies slightly beyond this bound; the difference
is consistent with rounding at the reference's displayed precision. Targeted repaired generations for
\texttt{batchdes}, \texttt{clay0204m}, and \texttt{risk2bpb} are retained as
regression evidence, separate from the uniform campaign.

This distinction is essential to interpreting the rejection counts: an exact
invalid inference refutes its submitted justification, whereas a failed
sufficient enclosure or curvature test may reflect the limits of the admitted
certificate rather than a false inequality or a nonconvex model.

\paragraph{Returned points and matching primal evidence.}
\label{sec:solver-audits}
Objective disagreements guide investigation; original-variable checks determine
what can be concluded about a returned point. For saved solver objectives
$r_{\rm solver}$, we include recorded GAMS model
statuses 1 (Optimal), 2 (Locally Optimal), and 8 (Integer Solution),\footnote{GAMS 53 User's Guide, \href{https://www.gams.com/53/docs/UG_GAMSOutput.html}{Model Status table}.}
and flag pairs satisfying
\[
 s(B-r_{\rm solver})>10^{-6}\max\{1,|B|\}.
\]
This one-sided criterion flags 18 historical bound/solver pairs. These are investigation candidates;
the two returned-point audits below support stronger, explicitly limited
conclusions.

For \texttt{clay0204m}, a saved SBB run reports objective 5075 with normal
completion and model status \emph{Integer Solution}; its log also warns of a
nonconvex model. We therefore do not describe it as an unqualified global
optimality status. Its printed point violates the two rows
\[
 x_2-x_4+51b_{13}\le46.5,\qquad
 x_6-x_7+86b_{24}\le82
\]
by 4.5 and 3, respectively. Allowing every printed coordinate an error of
$0.00005$ changes the row values by at most $0.00265$ and $0.0044$.
The violations are consequently robust to the displayed precision.
Independently, exact rational evaluation of a 52-variable point checks all
90 loaded constraints and 32 binary variables at objective 6545. A complete
certificate proves the matching lower bound 6545. Together these establish
an exact optimum for the loaded model. The archived point comes from rounded
BARON output with small distance coordinates explicitly replaced by zero;
its feasibility follows from the new rational evaluation, not from BARON's
status or tolerance.

For \texttt{risk2bpb}, the saved SHOT run reports an objective near
$-56.82087261$ and optimal model status. Its printed point assigns one to both
$b_4$ and $b_5$, although both variables are fixed at zero in the source model.
These unit violations suffice to reject that returned point. They do not
identify the internal cause of the discrepancy. The archived run used SHOT
1.1 with CPLEX 22.1; the baseline scripts requested 60 seconds, relative and
absolute optimality tolerances $10^{-4}$ and $10^{-6}$, and four threads.

To ensure that these statements compare the intended formulas, a separate
source audit matches the GAMS and loaded Python models under the same exact
binary64-constant interpretation. It eliminates the GAMS free objective
variable and defining equality, then compares each remaining constraint,
variable type, fixed value, and bound with an explicit index map. The
\texttt{clay0204m} comparison maps 53 variables and 91 rows to 52 variables
and 90 rows; \texttt{risk2bpb} maps 464 variables and 581 rows to 463 variables
and 580 rows. The audited violated rows and fixed bounds use exactly
representable coefficients, so their violations do not depend on a distinction
between decimal and binary64 constants. This is source-formula equivalence;
it does not verify the GAMS compiler or recover either solver's historical
in-memory model. The conclusions concern the saved returned points and the
specified source models, without attributing an unobserved internal defect.

The point audits establish concrete feasibility failures, and the separate
\texttt{clay0204m} witness establishes exact optimality for the loaded model.
Neither conclusion relies on treating a solver status or library objective as
a primal certificate. The source comparison states precisely which model is
being checked, while leaving the compiler and unseen solver state outside the
claim.

\subsection{Costs and practical limits}
\label{sec:experimental-costs}

Replay avoids repeating numerical search, but does not make large certificates
cheap to store or check. The uniform campaign takes 2,840.697 seconds of elapsed
generation time; its later separate replay takes 533.541 seconds. Accepted
checker calls sum to 2,791.064 seconds, with median 5.173 and maximum 221.484;
all 289 records sum to 2,953.982 seconds. These are observed shared-machine
costs with concurrent workers, not serial elapsed times or controlled speed
comparisons. Complete checking includes nonlinear/model checks, master matching,
and proof replay, rather than only discrete inference arithmetic. Historical
replay additionally includes external corroboration; its full costs are kept
separate in Appendix~\ref{app:experimental-accounting}.

The 203 accepted uniform proofs occupy 14,233,742,052 bytes and contain
11,749,855 derivations. The accepted historical proofs occupy
38,826,726,525 bytes; their median size is 74,469,242 bytes and their largest
is 2,163,285,624 bytes. Retaining failed completed proofs and fallback attempts
adds further audit value and storage cost. The complete proof collection is
30,664,561,063 bytes compressed and contains 93,713,729,595 regular-file bytes;
a small core supports representative replay without extracting that collection
(Appendix~\ref{app:reproduction}). These observations establish a practical
storage burden, not a certificate-size comparison with prior methods.

The machine provides approximately 30\,GiB of memory and 18 physical cores
(36 logical processors). Six instance workers share those resources. OA,
SCIP, and LP search request one thread each; proof completion retains its tool
defaults, so the entire pipeline is not single-threaded. The two-pass replay
in Section~\ref{sec:implementation} avoids retaining the whole proof text and
frees sparse rows after their last use. Its storage still includes per-derivation index arrays and depends on the
original master, input-line length, rational sizes, and simultaneously live
rows; machine capacity is not a measurement or bound on checker peak memory.
We infer no memory or speed advantage from these runs. The complete hardware,
software, timeout, phase-accounting, and hashing protocol is retained in
Appendix~\ref{app:experimental-accounting}.

A small portable core contains the checker sources, pinned dependency
specifications, all 289 campaign models, compact campaign records and exact comparison
tables, the source and returned-point audit, and representative certificates.
The separate bulk archive contains the completed historical and new proofs,
including rejected completed attempts. Generator dependencies are documented
separately: running the checker does not require Gurobi, IPOPT, SCIP, or GAMS.
License-identifying solver banners are removed only from distributed log
copies; the redaction manifest records original and distributed hashes.
Model files, mathematical data, and proof files are unmodified.

The practical result is a replayable collection with explicit limits: useful
bounds are available for many attempted models, exact audits can localize
specific failures, and expensive saved proofs remain separable from a small
checker-only starting point. The evidence supports those claims without
turning reference proximity, conservative rejection, or shared-machine timing
into stronger conclusions.
